\histitem{Operator algebras}

We have seen the role that Riesz, in expounding the work of the Hilbert school in 1913, assigned to the algebra \(\mathcal{L}(\mathrm{E})\) of endomorphisms of a Hilbert space, as well as to holomorphic functional calculus and the abstract notion of spectrum. In his 1916 memoir, he already seems to announce normed algebras as the natural framework for spectral theory; after the work of Banach and his school during the decade

of the 1920s, it is hardly surprising that this idea took a central place. The notion of Banach algebra was introduced by M. Nagumo {[}48{]} in 1936, while S. Mazur {[}44{]}, in 1938, showed that the only normed algebra over C which is a field is C itself (I, p. 26, cor. 2).

Toward the middle of the 1930s, in connection with spectral theory, opportunities multiplied for considering algebras of operators on Hilbert spaces abstractly. F. Murray and J. von Neumann, in a series of profound and influential works {[}47{]}, systematically studied the involutive subalgebras of \(\mathcal{L}(\mathrm{E})\) equal to their bicommutant (“von Neumann algebras”). On the other hand, in 1936, S. Steen proposed introducing an axiomatic notion of “algebra of operators” and studying the associated structures in depth {[}70{]}. Moreover, Stone {[}74{]}, motivated by the study of spectral projectors, studied in 1937 the unital algebras all of whose elements are idempotent (“Boolean algebras”: cf. ÉHM, p. 146). He noted that if X is a locally compact topological space, the characteristic functions of open sets and those of complements of everywhere dense sets generate, in \(\mathcal{C}(\mathrm{X})\) , a Boolean algebra A. He then proved that the ideals of A correspond naturally to the closed subsets of the space X, thereby introducing an idea that would prove particularly fruitful.

Beginning in 1939, Gelfand's viewpoint opened new perspectives on these relatively scattered results. His work seems to have been motivated in part by Wiener's ideas, in particular by his theorem on absolutely convergent Fourier series, as well as by the methods of Peter and Weyl. Between 1939 and 1946, in collaboration with M. Naimark, D. Raikov, and G. Shilov, Gelfand established the foundations of the theory of commutative Banach algebras and of the theory of stellar algebras (cf. {[}23{]}, vol. I, p. 169--400). In particular, he gave the classical proof of Wiener's theorem --- generalized by P. Lévy {[}39{]} --- which is found in this book (I, p. 38, example).

In Gelfand's approach, the essential object associated with a commutative Banach algebra A is the space X of maximal ideals of A. Indeed, via the Gelfand--Mazur theorem (I, p. 26, cor. 2), every element of A defines a function on X. It is well known how important this viewpoint would become in the subsequent evolution of algebraic geometry (cf. ÉHM, p. 146--148). This approach was itself motivated by Stone's work on Boolean algebras mentioned above.

It seems that this is L. Loomis, in a course published in 1953 {[}41, p. 53{]}, who first presented Gelfand theory with emphasis on the space of characters of A, defining the Gelfand transform in the sense in which we understand it in this book. The perhaps clearest advantage of this viewpoint is that the topological properties of the space of characters follow immediately from the compactness of the unit ball of the dual of a Banach space in the weak topology.

From his earliest works, Gelfand had constructed the holomorphic functional calculus in one variable in a Banach algebra by means of the Cauchy integral, and had deduced from it that an involutive Banach algebra admits a nontrivial decomposition as a product of rings when the space of its maximal ideals is not connected. The case of an arbitrary Banach algebra was treated only in 1953, when Shilov developed a form of the functional calculus in several variables. {[}68{]}.

As for star algebras, which Gelfand and Naimark had shown as early as 1942 to be realizable as algebras of operators on a Hilbert space (cf. V, p. 442, th. 2), they soon became the object of abundant and profound studies, spurred on by applications to group representations, by the work of Murray and von Neumann, and by the mathematical formalization of quantum mechanics. Let us mention in particular the contributions of I. Segal, who was one of the pioneers of the systematic study of representations of star algebras, in connection with the study of the unitary dual of locally compact groups, which we shall discuss in a few lines.

Although Gelfand, Naimark, Raikov, and Shilov had laid out the essential foundations of the theory of C*-algebras before 1946, certain technical points were improved later. Thus, R. Arens {[}1{]} showed how to avoid invoking the existence of the “Shilov boundary” (cf. I, p. 171, exerc. 26) in order to prove that the Gelfand transform is an involutive morphism. Moreover, the Gelfand–Naimark theory for unital C*-algebras initially added as an axiom the fact that \(1 + x^{*}x\) is invertible for all \(x\). They conjectured, however, that this axiom was unnecessary; the proof of this assertion, elementary but delicate, was only provided around 1952, following the work of M. Fukamiya {[}21{]} taken up by I. Kaplansky {[}62{]}.

